\documentclass[10pt,a4paper]{amsart}
\usepackage{amsmath,amssymb,amsthm,graphicx}
% Retained minimal preamble for the Stacks Project chapter varieties.tex.
% Verbatim excerpt of the chapter preamble: Xy-pic setup (native preamble lines 14-18), the native theorem
% declarations with the mathematical macros (native preamble lines 154-217), and the project own no-op
% definitions of the reference/slogan/history display environments (native wrapper.tex lines 1-3) that the
% counted statement uses for its slogan. Copied mechanically; nothing added, removed or redefined.
% For commutative diagrams we use Xy-pic
\usepackage[all]{xy}

% We use 2cell for 2-commutative diagrams.
\xyoption{2cell}
% Theorem environments.
%
\theoremstyle{plain}
\newtheorem{theorem}[subsection]{Theorem}
\newtheorem{proposition}[subsection]{Proposition}
\newtheorem{lemma}[subsection]{Lemma}

\theoremstyle{definition}
\newtheorem{definition}[subsection]{Definition}
\newtheorem{example}[subsection]{Example}
\newtheorem{exercise}[subsection]{Exercise}
\newtheorem{situation}[subsection]{Situation}

\theoremstyle{remark}
\newtheorem{remark}[subsection]{Remark}
\newtheorem{remarks}[subsection]{Remarks}

\numberwithin{equation}{subsection}

% Macros
%
\def\lim{\mathop{\mathrm{lim}}\nolimits}
\def\colim{\mathop{\mathrm{colim}}\nolimits}
\def\Spec{\mathop{\mathrm{Spec}}}
\def\Hom{\mathop{\mathrm{Hom}}\nolimits}
\def\Ext{\mathop{\mathrm{Ext}}\nolimits}
\def\SheafHom{\mathop{\mathcal{H}\!\mathit{om}}\nolimits}
\def\SheafExt{\mathop{\mathcal{E}\!\mathit{xt}}\nolimits}
\def\Sch{\mathit{Sch}}
\def\Mor{\mathop{\mathrm{Mor}}\nolimits}
\def\Ob{\mathop{\mathrm{Ob}}\nolimits}
\def\Sh{\mathop{\mathit{Sh}}\nolimits}
\def\NL{\mathop{N\!L}\nolimits}
\def\CH{\mathop{\mathrm{CH}}\nolimits}
\def\proetale{{pro\text{-}\acute{e}tale}}
\def\etale{{\acute{e}tale}}
\def\QCoh{\mathit{QCoh}}
\def\Ker{\mathop{\mathrm{Ker}}}
\def\Im{\mathop{\mathrm{Im}}}
\def\Coker{\mathop{\mathrm{Coker}}}
\def\Coim{\mathop{\mathrm{Coim}}}

% Boxtimes
%
\DeclareMathSymbol{\boxtimes}{\mathbin}{AMSa}{"02}

%
% Macros for moduli stacks/spaces
%
\def\QCohstack{\mathcal{QC}\!\mathit{oh}}
\def\Cohstack{\mathcal{C}\!\mathit{oh}}
\def\Spacesstack{\mathcal{S}\!\mathit{paces}}
\def\Quotfunctor{\mathrm{Quot}}
\def\Hilbfunctor{\mathrm{Hilb}}
\def\Curvesstack{\mathcal{C}\!\mathit{urves}}
\def\Polarizedstack{\mathcal{P}\!\mathit{olarized}}
\def\Complexesstack{\mathcal{C}\!\mathit{omplexes}}
% \Pic is the operator that assigns to X its picard group, usage \Pic(X)
% \Picardstack_{X/B} denotes the Picard stack of X over B
% \Picardfunctor_{X/B} denotes the Picard functor of X over B
\def\Pic{\mathop{\mathrm{Pic}}\nolimits}
\def\Picardstack{\mathcal{P}\!\mathit{ic}}
\def\Picardfunctor{\mathrm{Pic}}
\def\Deformationcategory{\mathcal{D}\!\mathit{ef}}
\newenvironment{reference}{}{}
\newenvironment{slogan}{}{}
\newenvironment{history}{}{}
\usepackage{cleveref}
\crefname{originallemma}{original lemma}{original lemmas}
\Crefname{originallemma}{Original lemma}{Original lemmas}
\makeatletter
\newlabel{algebra-lemma-CM-ring-catenary}{{00NM}{source}{}{sourcepointer-1}{}}
\newlabel{algebra-lemma-CM-ring-catenary@cref}{{[originallemma][0][]00NM}{[1][0][]source}{}{}{}}
\makeatother
\crefname{originallemma}{original lemma}{original lemmas}
\Crefname{originallemma}{Original lemma}{Original lemmas}
\makeatletter
\newlabel{algebra-lemma-catenary}{{02IH}{source}{}{sourcepointer-2}{}}
\newlabel{algebra-lemma-catenary@cref}{{[originallemma][0][]02IH}{[1][0][]source}{}{}{}}
\makeatother
\crefname{originallemma}{original lemma}{original lemmas}
\Crefname{originallemma}{Original lemma}{Original lemmas}
\makeatletter
\newlabel{algebra-lemma-dimension-at-a-point-finite-type-field}{{00P1}{source}{}{sourcepointer-3}{}}
\newlabel{algebra-lemma-dimension-at-a-point-finite-type-field@cref}{{[originallemma][0][]00P1}{[1][0][]source}{}{}{}}
\makeatother
\crefname{originallemma}{original lemma}{original lemmas}
\Crefname{originallemma}{Original lemma}{Original lemmas}
\makeatletter
\newlabel{algebra-lemma-dimension-at-a-point-finite-type-over-field}{{00OT}{source}{}{sourcepointer-4}{}}
\newlabel{algebra-lemma-dimension-at-a-point-finite-type-over-field@cref}{{[originallemma][0][]00OT}{[1][0][]source}{}{}{}}
\makeatother
\crefname{originallemma}{original lemma}{original lemmas}
\Crefname{originallemma}{Original lemma}{Original lemmas}
\makeatletter
\newlabel{algebra-lemma-dimension-closed-point-finite-type-field}{{00OU}{source}{}{sourcepointer-5}{}}
\newlabel{algebra-lemma-dimension-closed-point-finite-type-field@cref}{{[originallemma][0][]00OU}{[1][0][]source}{}{}{}}
\makeatother
\crefname{originallemma}{original lemma}{original lemmas}
\Crefname{originallemma}{Original lemma}{Original lemmas}
\makeatletter
\newlabel{algebra-lemma-dimension-height-polynomial-ring}{{00OR}{source}{}{sourcepointer-6}{}}
\newlabel{algebra-lemma-dimension-height-polynomial-ring@cref}{{[originallemma][0][]00OR}{[1][0][]source}{}{}{}}
\makeatother
\crefname{originallemma}{original lemma}{original lemmas}
\Crefname{originallemma}{Original lemma}{Original lemmas}
\makeatletter
\newlabel{algebra-lemma-dimension-prime-polynomial-ring}{{00P0}{source}{}{sourcepointer-7}{}}
\newlabel{algebra-lemma-dimension-prime-polynomial-ring@cref}{{[originallemma][0][]00P0}{[1][0][]source}{}{}{}}
\makeatother
\crefname{originallemma}{original lemma}{original lemmas}
\Crefname{originallemma}{Original lemma}{Original lemmas}
\makeatletter
\newlabel{algebra-lemma-dimension-spell-it-out}{{00OS}{source}{}{sourcepointer-8}{}}
\newlabel{algebra-lemma-dimension-spell-it-out@cref}{{[originallemma][0][]00OS}{[1][0][]source}{}{}{}}
\makeatother
\crefname{originallemma}{original lemma}{original lemmas}
\Crefname{originallemma}{Original lemma}{Original lemmas}
\makeatletter
\newlabel{algebra-lemma-quotient-catenary}{{00NK}{source}{}{sourcepointer-9}{}}
\newlabel{algebra-lemma-quotient-catenary@cref}{{[originallemma][0][]00NK}{[1][0][]source}{}{}{}}
\makeatother
\crefname{originallemma}{original lemma}{original lemmas}
\Crefname{originallemma}{Original lemma}{Original lemmas}
\makeatletter
\newlabel{lemma-locally-finite-type-Jacobson}{{0478}{source}{}{sourcepointer-10}{}}
\newlabel{lemma-locally-finite-type-Jacobson@cref}{{[originallemma][0][]0478}{[1][0][]source}{}{}{}}
\makeatother
\crefname{originaldefinition}{original definition}{original definitions}
\Crefname{originaldefinition}{Original definition}{Original definitions}
\makeatletter
\newlabel{topology-definition-catenary}{{02I1}{source}{}{sourcepointer-11}{}}
\newlabel{topology-definition-catenary@cref}{{[originaldefinition][0][]02I1}{[1][0][]source}{}{}{}}
\makeatother
\crefname{originaldefinition}{original definition}{original definitions}
\Crefname{originaldefinition}{Original definition}{Original definitions}
\makeatletter
\newlabel{topology-definition-equidimensional}{{0058}{source}{}{sourcepointer-12}{}}
\newlabel{topology-definition-equidimensional@cref}{{[originaldefinition][0][]0058}{[1][0][]source}{}{}{}}
\makeatother
\crefname{originallemma}{original lemma}{original lemmas}
\Crefname{originallemma}{Original lemma}{Original lemmas}
\makeatletter
\newlabel{topology-lemma-catenary}{{02I2}{source}{}{sourcepointer-13}{}}
\newlabel{topology-lemma-catenary@cref}{{[originallemma][0][]02I2}{[1][0][]source}{}{}{}}
\makeatother
\crefname{originallemma}{original lemma}{original lemmas}
\Crefname{originallemma}{Original lemma}{Original lemmas}
\makeatletter
\newlabel{topology-lemma-irreducible}{{004W}{source}{}{sourcepointer-14}{}}
\newlabel{topology-lemma-irreducible@cref}{{[originallemma][0][]004W}{[1][0][]source}{}{}{}}
\makeatother
\title{Dimension theory for locally algebraic k-schemes: catenarity, dimension at a point, and transcendence degree}
\author{The Stacks Project Authors}
\date{stacks-project:varieties.tex:local-snapshot-20261011; GNU FDL 1.2 or later}
\begin{document}
\maketitle
\noindent\textit{Editorial note.} All delivered mathematical text is the author's own, copied from the pinned Stacks Project chapter file varieties.tex; no statement or proof step was rewritten, reordered or removed.\par\medskip
\noindent\textit{Editorial note.} The author's opening paragraph explicitly announces the strategy of reducing to finite type k-algebras and citing the commutative algebra chapter; that paragraph is the first block of the proof and is reproduced as written.\par\medskip
\noindent\textit{Editorial note.} The counted spans carry 35 navigation stubs, one for each labelled item invoked outside the span; none of them replaces a load-bearing step, since the reductions themselves are written out inside the span. No bibliography citation occurs in the spans.\par\medskip
\noindent\textit{Editorial note.} The preamble is the chapter's retained minimal preamble (sources/stacks-varieties/retained-preamble.tex lines 1-77), a verbatim excerpt of the native preamble.\par\medskip
\noindent\textit{Editorial note.} The label lemma-dimension-locally-algebraic is defined only in varieties.tex, so no cross-chapter restatement had to be weighed for this card.\par\medskip
\noindent\textit{Allowed standard prerequisites (editorial).}\begin{itemize}
\item Catenary topological space. Topology, Definition topology-definition-catenary, named in the statement at line 3789. Supplies the notion verified in part (1).
\item Finitely generated k-algebras of finite type: catenarity, dimension at a point, dimension of polynomial rings. Algebra, Lemmas algebra-lemma-catenary, algebra-lemma-quotient-catenary, algebra-lemma-dimension-height-polynomial-ring, algebra-lemma-dimension-closed-point-finite-type-field, algebra-lemma-dimension-at-a-point-finite-type-field, algebra-lemma-dimension-at-a-point-finite-type-over-field, algebra-lemma-dimension-spell-it-out, algebra-lemma-dimension-prime-polynomial-ring. Supplies the affine case to which every part of the statement is reduced.
\item Normality criterion as the intersection of the localizations at height one primes. Algebra, Lemma algebra-lemma-normal-domain-intersection-localizations-height-1, used elsewhere in the chapter and cited in the sibling lemma; here the affine reductions rest on the corresponding algebra inputs. Supplies that affine pieces can be studied through their height one localizations.
\item Jacobson property of schemes locally of finite type over a field. Same chapter, Lemma lemma-locally-finite-type-Jacobson, invoked at proof line 3879. Supplies the closed point inside a nonempty open subset used in part (3).
\end{itemize}\medskip
\ifdefined\hypertarget\hypertarget{sourcepointer-1}{}\fi
\noindent\textit{Source reference.} Original lemma 00NM: sources/stacks-chow/tags.txt, source lines 864--864. Pointer outside the counted span, invoked once for context; the load-bearing argument of the counted statement is the author's proof reproduced in full in the proof span.\par\medskip
\ifdefined\hypertarget\hypertarget{sourcepointer-2}{}\fi
\noindent\textit{Source reference.} Original lemma 02IH: sources/stacks-chow/tags.txt, source lines 3300--3300. Pointer outside the counted span, invoked once for context; the load-bearing argument of the counted statement is the author's proof reproduced in full in the proof span.\par\medskip
\ifdefined\hypertarget\hypertarget{sourcepointer-3}{}\fi
\noindent\textit{Source reference.} Original lemma 00P1: sources/stacks-chow/tags.txt, source lines 915--915. Pointer outside the counted span, invoked once for context; the load-bearing argument of the counted statement is the author's proof reproduced in full in the proof span.\par\medskip
\ifdefined\hypertarget\hypertarget{sourcepointer-4}{}\fi
\noindent\textit{Source reference.} Original lemma 00OT: sources/stacks-chow/tags.txt, source lines 907--907. Pointer outside the counted span, invoked once for context; the load-bearing argument of the counted statement is the author's proof reproduced in full in the proof span.\par\medskip
\ifdefined\hypertarget\hypertarget{sourcepointer-5}{}\fi
\noindent\textit{Source reference.} Original lemma 00OU: sources/stacks-chow/tags.txt, source lines 908--908. Pointer outside the counted span, invoked once for context; the load-bearing argument of the counted statement is the author's proof reproduced in full in the proof span.\par\medskip
\ifdefined\hypertarget\hypertarget{sourcepointer-6}{}\fi
\noindent\textit{Source reference.} Original lemma 00OR: sources/stacks-chow/tags.txt, source lines 905--905. Pointer outside the counted span, invoked once for context; the load-bearing argument of the counted statement is the author's proof reproduced in full in the proof span.\par\medskip
\ifdefined\hypertarget\hypertarget{sourcepointer-7}{}\fi
\noindent\textit{Source reference.} Original lemma 00P0: sources/stacks-chow/tags.txt, source lines 914--914. Pointer outside the counted span, invoked once for context; the load-bearing argument of the counted statement is the author's proof reproduced in full in the proof span.\par\medskip
\ifdefined\hypertarget\hypertarget{sourcepointer-8}{}\fi
\noindent\textit{Source reference.} Original lemma 00OS: sources/stacks-chow/tags.txt, source lines 906--906. Pointer outside the counted span, invoked once for context; the load-bearing argument of the counted statement is the author's proof reproduced in full in the proof span.\par\medskip
\ifdefined\hypertarget\hypertarget{sourcepointer-9}{}\fi
\noindent\textit{Source reference.} Original lemma 00NK: sources/stacks-chow/tags.txt, source lines 862--862. Pointer outside the counted span, invoked once for context; the load-bearing argument of the counted statement is the author's proof reproduced in full in the proof span.\par\medskip
\ifdefined\hypertarget\hypertarget{sourcepointer-10}{}\fi
\noindent\textit{Source reference.} Original lemma 0478: sources/stacks-varieties/pinned-original-source.tex, source lines 2670--2670. Pointer outside the counted span, invoked once for context; the load-bearing argument of the counted statement is the author's proof reproduced in full in the proof span.\par\medskip
\ifdefined\hypertarget\hypertarget{sourcepointer-11}{}\fi
\noindent\textit{Source reference.} Original definition 02I1: sources/stacks-chow/tags.txt, source lines 3284--3284. Pointer outside the counted span, invoked once for context; the load-bearing argument of the counted statement is the author's proof reproduced in full in the proof span.\par\medskip
\ifdefined\hypertarget\hypertarget{sourcepointer-12}{}\fi
\noindent\textit{Source reference.} Original definition 0058: sources/stacks-chow/tags.txt, source lines 200--200. Pointer outside the counted span, invoked once for context; the load-bearing argument of the counted statement is the author's proof reproduced in full in the proof span.\par\medskip
\ifdefined\hypertarget\hypertarget{sourcepointer-13}{}\fi
\noindent\textit{Source reference.} Original lemma 02I2: sources/stacks-chow/tags.txt, source lines 3285--3285. Pointer outside the counted span, invoked once for context; the load-bearing argument of the counted statement is the author's proof reproduced in full in the proof span.\par\medskip
\ifdefined\hypertarget\hypertarget{sourcepointer-14}{}\fi
\noindent\textit{Source reference.} Original lemma 004W: sources/stacks-chow/tags.txt, source lines 188--188. Pointer outside the counted span, invoked once for context; the load-bearing argument of the counted statement is the author's proof reproduced in full in the proof span.\par\medskip
\begin{lemma}
\label{lemma-dimension-locally-algebraic}
Let $k$ be a field. Let $X$ be a locally algebraic $k$-scheme.
\begin{enumerate}
\item
\label{item-catenary}
The topological space of $X$ is catenary
(Topology, Definition \ref{topology-definition-catenary}).
\item
\label{item-dimension-at-closed-point}
For $x \in X$ closed, we have $\dim_x(X) = \dim(\mathcal{O}_{X, x})$.
\item
\label{item-dimension-irreducible}
For $X$ irreducible we have $\dim(X) = \dim(U)$ for any
nonempty open $U \subset X$ and $\dim(X) = \dim_x(X)$
for any $x \in X$.
\item
\label{item-irreducible-maximal-chains}
For $X$ irreducible any chain of irreducible closed subsets can be
extended to a maximal chain and all maximal chains of irreducible
closed subsets have length equal to $\dim(X)$.
\item
\label{item-dimension-irreducibles-passing-through}
For $x \in X$ we have
$\dim_x(X) = \max \dim(Z) = \min \dim(\mathcal{O}_{X, x'})$
where the maximum is over irreducible components
$Z \subset X$ containing $x$ and the minimum is over
specializations $x \leadsto x'$ with $x'$ closed in $X$.
\item
\label{item-dimension-irreducible-trdeg}
If $X$ is irreducible with generic point $x$, then
$\dim(X) = \text{trdeg}_k(\kappa(x))$.
\item
\label{item-immediate-specialization}
If $x \leadsto x'$ is an immediate specialization
of points of $X$, then we have
$\text{trdeg}_k(\kappa(x)) = \text{trdeg}_k(\kappa(x')) + 1$.
\item
\label{item-dimension-sup-trdeg}
The dimension of $X$ is the supremum of the numbers
$\text{trdeg}_k(\kappa(x))$ where $x$ runs over the
generic points of the irreducible components of $X$.
\item
\label{item-specialization}
If $x \leadsto x'$ is a nontrivial specialization of points of $X$, then
\begin{enumerate}
\item $\dim_x(X) \leq \dim_{x'}(X)$,
\item $\dim(\mathcal{O}_{X, x}) < \dim(\mathcal{O}_{X, x'})$,
\item $\text{trdeg}_k(\kappa(x)) > \text{trdeg}_k(\kappa(x'))$, and
\item any maximal chain of nontrivial specializations
$x = x_0 \leadsto x_1 \leadsto \ldots \leadsto x_n = x'$ has
length $n = \text{trdeg}_k(\kappa(x)) - \text{trdeg}_k(\kappa(x'))$.
\end{enumerate}
\item
\label{item-dimension-formula}
For $x \in X$ we have
$\dim_x(X) = \text{trdeg}_k(\kappa(x)) + \dim(\mathcal{O}_{X, x})$.
\item
\label{item-immediate-specialization-local-ring}
If $x \leadsto x'$ is an immediate specialization
of points of $X$ and $X$ is irreducible or equidimensional, then
$\dim(\mathcal{O}_{X, x'}) = \dim(\mathcal{O}_{X, x}) + 1$.
\end{enumerate}
\end{lemma}

\begin{proof}
Instead on relying on the more general results proved earlier
we will reduce the statements to the corresponding statements
for finite type $k$-algebras and cite results from the chapter
on commutative algebra.

\medskip\noindent
Proof of (\ref{item-catenary}). This is local on $X$ by
Topology, Lemma \ref{topology-lemma-catenary}. Thus we may
assume $X = \Spec(A)$ where $A$ is a finite type $k$-algebra.
We have to show that $A$ is catenary
(Algebra, Lemma \ref{algebra-lemma-catenary}).
We can reduce to $k[x_1, \ldots, x_n]$ using
Algebra, Lemma \ref{algebra-lemma-quotient-catenary} and then apply
Algebra, Lemma \ref{algebra-lemma-dimension-height-polynomial-ring}.
Alternatively, this holds because $k$ is Cohen-Macaulay (trivially) and
Cohen-Macaulay rings are universally catenary
(Algebra, Lemma \ref{algebra-lemma-CM-ring-catenary}).

\medskip\noindent
Proof of (\ref{item-dimension-at-closed-point}). Choose an affine
neighbourhood $U = \Spec(A)$ of $x$. Then $\dim_x(X) = \dim_x(U)$.
Hence we reduce to the affine case, which is
Algebra, Lemma \ref{algebra-lemma-dimension-closed-point-finite-type-field}.

\medskip\noindent
Proof of (\ref{item-dimension-irreducible}). It suffices to show that
any two nonempty affine opens $U, U' \subset X$ have the same dimension
(any finite chain of irreducible subsets meets an affine open).
Pick a closed point $x$ of $X$ with $x \in U \cap U'$. This is possible
because $X$ is irreducible, hence $U \cap U'$ is nonempty, hence there is
such a closed point because $X$ is Jacobson by
Lemma \ref{lemma-locally-finite-type-Jacobson}. Then
$\dim(U) = \dim(\mathcal{O}_{X, x}) = \dim(U')$ by
Algebra, Lemma \ref{algebra-lemma-dimension-spell-it-out}
(strictly speaking you have to replace $X$ by its reduction before
applying the lemma).

\medskip\noindent
Proof of (\ref{item-irreducible-maximal-chains}). Given a chain
of irreducible closed subsets we can find an affine open $U \subset X$
which meets the smallest one. Thus the statement follows from
Algebra, Lemma \ref{algebra-lemma-dimension-spell-it-out} 
and $\dim(U) = \dim(X)$ which we have seen in
(\ref{item-dimension-irreducible}).

\medskip\noindent
Proof of (\ref{item-dimension-irreducibles-passing-through}).
Choose an affine neighbourhood $U = \Spec(A)$ of $x$. Then
$\dim_x(X) = \dim_x(U)$. The rule $Z \mapsto Z \cap U$
is a bijection between irreducible components of $X$ passing through
$x$ and irreducible components of $U$ passing through $x$.
Also, $\dim(Z \cap U) = \dim(Z)$ for such $Z$ by
(\ref{item-dimension-irreducible}).
Hence the statement follows from
Algebra, Lemma \ref{algebra-lemma-dimension-at-a-point-finite-type-over-field}.

\medskip\noindent
Proof of (\ref{item-dimension-irreducible-trdeg}). By
(\ref{item-dimension-irreducible}) this reduces to the case where
$X = \Spec(A)$ is affine. In this case it follows from
Algebra, Lemma \ref{algebra-lemma-dimension-prime-polynomial-ring}
applied to $A_{red}$.

\medskip\noindent
Proof of (\ref{item-immediate-specialization}).
Let $Z = \overline{\{x\}} \supset Z' = \overline{\{x'\}}$.
Then it follows from (\ref{item-irreducible-maximal-chains}) that
$Z \supset Z'$ is the start of a maximal chain of
irreducible closed subschemes in $Z$
and consequently $\dim(Z) = \dim(Z') + 1$.
We conclude by (\ref{item-dimension-irreducible-trdeg}).

\medskip\noindent
Proof of (\ref{item-dimension-sup-trdeg}). A simple topological argument
shows that $\dim(X) = \sup \dim(Z)$ where the supremum is over the
irreducible components of $X$ (hint: use
Topology, Lemma \ref{topology-lemma-irreducible}).
Thus this follows from (\ref{item-dimension-irreducible-trdeg}).

\medskip\noindent
Proof of (\ref{item-specialization}). Part (a) follows from the
fact that any open $U \subset X$ containing $x'$ also contains $x$.
Part (b) follows because $\mathcal{O}_{X, x}$ is a localization of
$\mathcal{O}_{X, x'}$ hence any chain of primes in $\mathcal{O}_{X, x}$
corresponds to a chain of primes in $\mathcal{O}_{X, x'}$ which can
be extended by adding $\mathfrak m_{x'}$ at the end.
Both (c) and (d) follow formally from (\ref{item-immediate-specialization}).

\medskip\noindent
Proof of (\ref{item-dimension-formula}). Choose an affine
neighbourhood $U = \Spec(A)$ of $x$. Then $\dim_x(X) = \dim_x(U)$.
Hence we reduce to the affine case, which is
Algebra, Lemma \ref{algebra-lemma-dimension-at-a-point-finite-type-field}.

\medskip\noindent
Proof of (\ref{item-immediate-specialization-local-ring}).
If $X$ is equidimensional
(Topology, Definition \ref{topology-definition-equidimensional})
then $\dim(X)$ is equal to the dimension of every irreducible
component of $X$, whence $\dim_x(X) = \dim(X) = \dim_{x'}(X)$
by (\ref{item-dimension-irreducibles-passing-through}).
Thus this follows from (\ref{item-immediate-specialization}).
\end{proof}

\end{document}
